\documentclass[12pt]{article}
\def\activityid{F08-M-v2}\usepackage[letterpaper,margin=.65in,top=.60in,bottom=.65in]{geometry}
\usepackage[T1]{fontenc}
\usepackage[default]{sourcesanspro}
\usepackage{microtype,tikz,array,tabularx,fancyhdr,amsmath,amssymb}
\usepackage[hidelinks]{hyperref}
\usetikzlibrary{calc,arrows.meta}
\setlength{\parindent}{0pt}\setlength{\parskip}{7pt}
\setlength{\headheight}{0pt}\setlength{\footskip}{26pt}\setlength{\emergencystretch}{2em}
\pagestyle{fancy}\fancyhf{}\renewcommand{\headrulewidth}{0pt}\renewcommand{\footrulewidth}{.4pt}
\fancyfoot[L]{\scriptsize Bellingham Math Circle / Week 8 / \activityid}
\fancyfoot[R]{\scriptsize \thepage}
\definecolor{ink}{gray}{.12}\definecolor{soft}{gray}{.95}\color{ink}
\newcommand{\PageTitle}[2]{{\fontsize{10}{12}\selectfont\bfseries WEEK 8\quad STRATEGY LAB\hfill #1\par}\vspace{3pt}{\fontsize{26}{29}\selectfont\bfseries #2\par}\vspace{2pt}}
\newcommand{\NameLine}{{\small Name: \rule{2.65in}{.4pt}\hfill Date: \rule{1.35in}{.4pt}\par}\vspace{4pt}}
\newcommand{\Task}[2]{\par\vspace{3pt}{\large\bfseries #1}\par #2\par}
\newcommand{\Subhead}[1]{\par\vspace{3pt}{\large\bfseries #1}\par}
\newcommand{\RuleBox}[1]{\par\begingroup\setlength{\fboxsep}{8pt}\colorbox{soft}{\parbox{\dimexpr\linewidth-16pt\relax}{#1}}\endgroup\par}
\newcommand{\WriteBox}[2]{\par\textbf{#1}\par\vspace{2pt}\begin{tikzpicture}\draw[black!40,line width=.5pt] (0,0) rectangle (\dimexpr\linewidth-.5pt\relax,#2);\end{tikzpicture}\par}
\newcommand{\StudentFooter}{\vfill{\small Try it with objects. Explain by pointing, talking, or drawing.}\par}
\newcommand{\RookBoard}[3]{\begin{tikzpicture}[x=#3,y=#3]\draw[step=1,black!45] (0,0) grid (#1,#2);\draw[thick] (0,0) rectangle (#1,#2);\node[font=\Large] at ({#1-.5},.5){$\star$};\draw[-{Latex},thick] (0,-.35)--(#1,-.35);\draw[-{Latex},thick] ({#1+.35},#2)--({#1+.35},0);\end{tikzpicture}}
\newcommand{\Table}[3]{\begin{tikzpicture}[x=#3,y=#3]\draw[step=1,black!25] (0,0) grid (#1,#2);\draw[thick] (0,0) rectangle (#1,#2);\fill (0,0)circle(3pt);\draw[-{Latex},very thick] (0,0)--(.7,.7);\node[below] at ({#1/2},-.1){width #1};\node[rotate=90,above] at (-.2,{#2/2}){height #2};\end{tikzpicture}}
\newcommand{\GraphNodes}{\coordinate(A)at(0,0);\coordinate(B)at(3,0);\coordinate(C)at(1.5,2.6);\coordinate(D)at(1.5,.95);}
\newcommand{\Kfour}[1]{\begin{tikzpicture}[scale=#1]\GraphNodes\draw[very thick](A)--(B)--(C)--(A)--(D)--(B) (D)--(C);\foreach \n in {A,B,C,D}{\node[circle,draw,fill=white,minimum size=22pt,font=\bfseries]at(\n){\n};}\end{tikzpicture}}
\newcommand{\House}[1]{\begin{tikzpicture}[scale=#1]\coordinate(A)at(0,0);\coordinate(B)at(3,0);\coordinate(C)at(3,2);\coordinate(D)at(0,2);\coordinate(E)at(1.5,3.3);\draw[very thick](A)--(B)--(C)--(D)--(A) (D)--(E)--(C);\foreach \n in {A,B,C,D,E}{\node[circle,draw,fill=white,minimum size=22pt,font=\bfseries]at(\n){\n};}\end{tikzpicture}}

\begin{document}

\PageTitle{GRADES 2--3}{Map the traps}\NameLine
\RuleBox{Share one token. Move right OR down any positive number of squares. No diagonal moves. Reach the star to win.}
\Task{1. Play, then make a map.}{Try different starts. Mark a square T if the player who receives it must lose against perfect replies. Mark W if that player can force a win. Use ? while unsure.}
\begin{center}\RookBoard{4}{4}{.65in}\end{center}
Treat the star as already finished: the next player has no move. Start your map near it. A single lost game does not prove T!
\Task{2. A claim that survives every reply.}{Choose one T. Show that \textbf{every} move from it gives the next player a W. Then choose one W and show \textbf{one} move to a T.}
\WriteBox{My two squares and reasons:}{1.15in}
\StudentFooter

\newpage

\PageTitle{GRADES 2--3}{The board is two piles}\NameLine
\Task{3. Count what is left.}{Put the token 3 squares left and 2 squares above the star. Make a pile of 3 and a pile of 2. A right move shrinks the first pile; a down move shrinks the second.}
\RuleBox{Pile rule: remove any positive number from ONE pile. Take the last counter to win. Play both versions and match their moves.}
\Task{4. Prove a rule for every rectangle.}{Which two-pile positions are traps? Explain these two things with counters:}
\begin{enumerate}\item Every move from a trap breaks your pattern.\item From any other position, you can make the pattern.\end{enumerate}
\WriteBox{My pattern and why both statements hold:}{1.1in}
\Task{5. Add a third pile.}{Keep the same pile rule. Try \((1,1,1)\), then \((1,2,3)\). Does ``two matching piles'' still guarantee a trap? List all first moves from \((1,2,3)\) and find a winning reply to each.}
\WriteBox{A surprise or a new conjecture:}{.85in}
\textbf{Make a puzzle:} Choose a two-pile start with exactly one winning move. Let a partner test it.
\StudentFooter

\end{document}
